The Modifiable Areal Unit Problem (MAUP), first systematically described by
\citet{openshaw1984modifiable}, is one of the foundational challenges of
quantitative geography: the same continuous spatial phenomenon produces
different numerical summaries depending on how analysts divide the territory
into discrete units.  MAUP has two components.  The \emph{scale effect}
occurs when aggregating from fine to coarse resolution changes summary
statistics even when no underlying data are altered.  The \emph{zoning
effect} occurs when different boundary configurations at the same resolution
produce different results.  Both effects are pervasive in spatial social
science \citep{openshaw1983modifiable, jelinski1996modifiable}.

Congressional redistricting is acutely vulnerable to MAUP.  Every
algorithmic redistricting method treats geographic space as a discrete graph
whose nodes are some census unit---blocks, block groups, tracts, or county
equivalents.  The choice of unit is not prescribed by any constitutional or
statutory requirement.  Researchers have overwhelmingly settled on census
tracts as a pragmatic default: tracts are small enough for fine-grained
boundary placement, large enough to have stable demographic estimates, and
small enough that national graphs are computationally tractable
\citep{census_tracts_reference}.  Nevertheless, state redistricting
commissions routinely use census blocks for official plans
\citep{ncsl_redistricting}, creating a potential inconsistency between
research findings and operational practice.

This inconsistency motivates a direct empirical test.  If key findings from
algorithmic redistricting---compactness scores, partisan seat predictions,
minority representation counts---change materially when the spatial unit
changes from tracts to block groups or county equivalents, then researchers
who use tracts cannot confidently generalise to practice.  If findings are
robust across the unit-count range, then the tract-level literature stands on
firm methodological ground.

\subsection{Why Bisect Is a Revealing Test Case}

The recursive bisection algorithm used in this study optimises
\emph{graph-edge cuts}: it partitions a planar graph so as to minimise the
total length-weighted boundary between districts.  Crucially, the objective
function does not reference polygon area or perimeter directly---it operates
on the abstract topology of the adjacency graph.  This scale-free structure
suggests a hypothesis: \emph{bisect should be comparatively MAUP-robust
because it never assumes anything about the physical size of nodes.}

Testing this hypothesis requires varying the unit while holding the algorithm
and parameters constant.  We do exactly that, running identical bisect
configurations at three spatial resolutions spanning the census hierarchy.

\subsection{Unit-Count Range}

The three resolution levels tested span a 130$\times$ unit-count range
nationally:

\begin{itemize}
  \item \textbf{County equivalents}: $\approx 3{,}100$ units nationally;
        mean population $\approx 100{,}000$.  Used as the \emph{coarse}
        baseline.
  \item \textbf{Census tracts}: $\approx 85{,}000$ units nationally;
        mean population $\approx 4{,}000$.  Used as the \emph{standard}
        baseline, consistent with the broader bisect portfolio.
  \item \textbf{Block groups}: $\approx 240{,}000$ units nationally;
        mean population $\approx 1{,}200$.  Used as the \emph{fine}
        resolution test.
\end{itemize}

The ratio between the largest (county-equivalent, $\sim$100,000 pop) and
smallest (block group, $\sim$1,200 pop) units is approximately $83\times$ in
population terms and $240{,}000 / 3{,}100 \approx 77\times$ in unit count.
Including the full tract-to-county range yields the headline $130\times$
figure (block groups to county equivalents by unit count, $240k$ vs $3k$).

\subsection{Contributions}

This paper makes three contributions.  First, it provides the first
systematic MAUP analysis in algorithmic redistricting, testing an untested
but important methodological assumption.  Second, it identifies the
structural mechanism---graph-distance-based optimisation---that makes bisect
comparatively robust.  Third, it establishes bounds: the two outlier states
(Idaho and Wyoming) exhibit $>2\%$ compactness variation, driven by their
extremely small district counts ($k = 2$) rather than by MAUP per se.

The remainder of the paper is organised as follows.  Section~2 reviews MAUP
theory, census geographic hierarchy, and the existing redistricting
literature's treatment of spatial resolution.  Section~3 describes the
experimental design.  Section~4 presents results across five focus states and
the full 50-state summary.  Section~5 discusses the mechanism and
implications.  Section~6 concludes.
